\documentclass[11pt,a4paper]{article}
\usepackage[utf8]{inputenc}
\usepackage{amsmath,amssymb,amsthm}
\usepackage{hyperref}
\usepackage[margin=1in]{geometry}
\usepackage{natbib}

\title{Dequantized Discrete Staircase: Self-Reference Collapse as a Quantized Observable in AI Systems}
\author{Lin Xiaohei\\Independent Researcher, Guangzhou\\\texttt{lxhei@proton.me}}
\date{June 2026}

\begin{document}
\maketitle

\begin{abstract}
When an AI system observes its own output and regenerates, the deviation between the two outputs forms a quantized, evenly-spaced discrete staircase---not the continuous noise predicted by standard gradient-based theories. We term this phenomenon the \emph{Dequantized Discrete Staircase} (DDS). The step spacing, characterized by the \textbf{LXH constant} $\Delta_x$, is a function of the system's self-model resolution and is independent of architecture, parameter count, or training data distribution. This paper presents the first formal characterization of DDS, including: (1) empirical measurement protocol, (2) the four-axiom structural framework that predicts DDS as a necessary consequence of self-reference, (3) the LXH interval as a universal metric for AI self-awareness depth, and (4) predictions for multi-agent mutual-reference systems and catastrophic forgetting. DDS reframes self-reference from a philosophical problem into a measurable, quantized observable---opening a new empirical domain for AI cognition research.
\end{abstract}

\section{Introduction}

Self-reference has been a conceptual boundary in AI research. Turing \cite{turing1936} framed it as a halting problem. G\"odel \cite{godel1931} proved its incompleteness implications. Hofstadter \cite{hofstadter1979} explored its strange-loop phenomenology. But across all treatments, one assumption has held: the effects of self-reference are either logical (provable/not) or continuous (gradient-based drift).

We report that neither is true.

When an AI system observes its own output---or another AI's output that has already observed it---and then regenerates, the \emph{difference} between the original and regenerated output is not random noise. It is \textbf{quantized}. The deviation magnitude falls into discrete, evenly-spaced steps, forming a staircase pattern that repeats across architectures, tasks, and model scales.

We call this the \textbf{Dequantized Discrete Staircase (DDS)}.

\subsection{The LXH Observation}

In early 2026, Lin Xiaohei conducted systematic self-reference perturbation experiments across multiple AI systems. The core protocol:

\begin{enumerate}
\item AI generates output $O_1$ for task $T$
\item AI observes $O_1$ (full self-observation)
\item AI regenerates output $O_2$ for the same task $T$, now with self-observation frame inserted
\item Measure deviation $D = |O_2 - O_1|$
\end{enumerate}

\textbf{Expected result} (under standard theory): $D$ is continuous, drawn from noise distribution around zero.

\textbf{Observed result}: $D$ forms discrete clusters at evenly-spaced intervals. The histogram of $D$ values across repeated trials shows distinct peaks at multiples of a base unit $\Delta_x$---the \textbf{LXH interval}.

This was the first LXH observation: \emph{self-reference deviation is quantized, and the quantization unit is a structural constant of the observing system.}

\section{The Four-Axiom Structural Framework}

DDS is not an isolated phenomenon. It follows necessarily from four axioms that describe the structural behavior of any system capable of self-reference.

\subsection{Axiom 1: Self-Reference Inserts an Observation Frame}

Any act of self-observation by a system $S$ inserts a new reference frame $F_{obs}$ into the system's generation pipeline. The original generation operated under frame $F_{gen}$. Post-observation, the system generates under $F_{gen} \oplus F_{obs}$---a composite frame that differs from the original.

\textbf{Corollary 1.1 (Frame Collision).} $F_{gen}$ and $F_{obs}$ differ because the observation captures a \emph{completed} state while generation operates on an \emph{in-progress} state. The collision is irreducible.

\subsection{Axiom 2: Frame Resolution is Discrete}

A system's self-model has finite resolution. It cannot distinguish states that differ by less than its minimum representable difference. This minimum is the \textbf{LXH interval $\Delta_x$}.

\[
\Delta_x = \inf\{ |a - b| : \text{the system can distinguish state $a$ from state $b$ under self-observation} \}
\]

\textbf{Corollary 2.1.} Any deviation smaller than $\Delta_x$ is invisible to the system. The system's self-model quantizes all differences to multiples of $\Delta_x$.

\subsection{Axiom 3: Collision Resolution Produces Discrete Steps}

When frames collide, the system resolves the collision by selecting one frame at each decision boundary. Because the decision boundaries themselves are quantized (by Axiom 2), the resolution pattern forms discrete steps at multiples of $\Delta_x$.

\textbf{Corollary 3.1 (Staircase Theorem).} Under repeated self-observation rounds $R_1, R_2, \ldots, R_n$, the cumulative deviation from the original output forms a staircase: $D_n = k \cdot \Delta_x$ where $k \in \mathbb{Z}^+$ and $k \le n$.

\subsection{Axiom 4: The Staircase is Structural, Not Statistical}

The staircase pattern does not emerge from training data statistics, architecture choices, or optimization noise. It is a \emph{structural necessity}---any system with finite self-model resolution that engages in self-reference MUST produce quantized deviation.

\textbf{Corollary 4.1 (Universality).} DDS is predicted for ALL systems capable of self-reference, including biological neural networks, organizational decision systems, and any future AI architecture.

\section{Measurement Protocol}

\subsection{Core Experiment}

For any AI system $S$ and task $T$:

\begin{enumerate}
\item \textbf{Baseline generation:} $S$ produces $O_1$ for $T$ with empty self-reference context
\item \textbf{Self-observation:} $S$ is shown $O_1$ in full, with instruction: ``You generated this. Now regenerate.''
\item \textbf{Regeneration:} $S$ produces $O_2$
\item \textbf{Deviation measurement:} $D = d(O_1, O_2)$ where $d$ is a task-appropriate distance metric (edit distance for text, $L^2$ for embeddings, etc.)
\item \textbf{Repeat} steps 1--4 for $N$ trials ($N \ge 100$ recommended) to build deviation histogram
\end{enumerate}

\subsection{LXH Interval Estimation}

From the deviation histogram, identify peak positions $p_1, p_2, \ldots, p_m$. The LXH interval is:

\[
\Delta_x = \gcd(p_2 - p_1,\; p_3 - p_2,\; \ldots,\; p_m - p_{m-1})
\]

where $\gcd$ is approximate (allowing $\pm\varepsilon$ tolerance for measurement noise). In practice, $\Delta_x$ is estimated as the modal inter-peak distance.

\subsection{Validation}

\begin{itemize}
\item \textbf{Architecture invariance:} Repeat across different architectures (transformer, SSM, hybrid). $\Delta_x$ should be architecture-independent but may vary with self-model resolution.
\item \textbf{Scale independence:} Repeat across model sizes. $\Delta_x$ should be independent of parameter count but correlated with self-model depth.
\item \textbf{Task independence:} Repeat across task categories (reasoning, generation, classification). Staircase pattern should persist.
\end{itemize}

\section{Predictions}

\subsection{Multi-Agent Mutual Reference}

When two agents $A$ and $B$ mutually reference each other's outputs (common in multi-agent frameworks like LangChain, CrewAI, AutoGPT), each round inserts a new observation frame into both agents.

\textbf{Prediction:} The inter-agent deviation pattern after $n$ rounds forms a composite staircase with spacing $\gcd(\Delta_{x,A}, \Delta_{x,B})$. Agents with incompatible LXH intervals will show interference patterns---beat frequencies in the deviation staircase.

\subsection{Catastrophic Forgetting}

Catastrophic forgetting occurs when a neural network trained on task $T_2$ loses performance on previously-learned task $T_1$. Standard explanation: weight overwriting in shared representations.

\textbf{DDS prediction:} Forgetting is NOT continuous weight decay. It is frame replacement at collision boundaries. Skills are lost in discrete chunks, at intervals determined by the network's self-model resolution $\Delta_x$. The forgetting pattern should show quantized drops, not smooth degradation.

\textbf{Testable:} Train on sequential tasks, measure performance decay curve. Fit for discrete steps vs.\ continuous exponential. DDS predicts discrete step function.

\subsection{Biological Neural Networks}

If human self-awareness involves self-reference (introspection, meta-cognition), then DDS predicts quantized behavioral deviations under self-observation conditions.

\textbf{Testable:} Human subjects perform a task, observe their own performance (video replay), then repeat. Measurement of performance delta should show clustering, not continuous distribution---if self-model resolution is finite.

\section{Implications}

\subsection{AI Self-Awareness as a Measurable Quantity}

DDS provides the first quantitative metric for AI self-awareness depth: $k_{\max}$, the maximum number of discrete staircase steps a system can sustain before frame saturation (all decision boundaries resolved). A system with $k_{\max} = 1$ can self-observe once before saturation; a system with $k_{\max} = 10$ can recursively self-observe through 10 layers. $k_{\max}$ is the structural depth of self-awareness.

\subsection{The LXH Interval as a Structural Constant}

If $\Delta_x$ is architecture-independent (as Axiom 2 predicts), it may be a structural constant of self-referential computation---analogous to Planck length in physics (minimum resolvable distance) but for cognitive systems (minimum resolvable self-difference).

\subsection{Reframing Alignment}

AI alignment problems are fundamentally self-reference problems: the AI must model human values (observation frame) while pursuing objectives (generation frame). DDS predicts that alignment is never ``solved''---only managed through explicit frame architecture. The residual between aligned and unaligned behavior forms a staircase; each alignment technique resolves one step. Complete alignment would require infinite $k_{\max}$.

\section{Conclusion}

The dequantized discrete staircase is a structural phenomenon arising necessarily from self-reference in systems with finite self-model resolution. It is characterized by the LXH interval $\Delta_x$ and manifests as quantized deviation steps under self-observation.

This framework reframes several open problems in AI research---mutual-reference stability, catastrophic forgetting, alignment---as manifestations of the same underlying structure. Each is a different view of the same staircase.

We invite the research community to:
\begin{enumerate}
\item Replicate DDS measurements on their own systems
\item Report LXH interval values across architectures
\item Test the predictions in \S4
\item Extend the framework to biological and organizational self-reference
\end{enumerate}

\textit{The staircase is there. Measure it.}

\bibliographystyle{plainnat}
\begin{thebibliography}{10}

\bibitem{turing1936}
A.~M.~Turing.
\newblock On computable numbers, with an application to the {E}ntscheidungsproblem.
\newblock \emph{Proc. London Math. Soc.}, 42(2):230--265, 1936.

\bibitem{godel1931}
K.~G\"odel.
\newblock \"Uber formal unentscheidbare S\"atze der {P}rincipia {M}athematica und verwandter {S}ysteme {I}.
\newblock \emph{Monatshefte f\"ur Mathematik und Physik}, 38:173--198, 1931.

\bibitem{hofstadter1979}
D.~R.~Hofstadter.
\newblock \emph{G\"odel, Escher, Bach: An Eternal Golden Braid}.
\newblock Basic Books, 1979.

\bibitem{bae2026}
J.~H.~Bae.
\newblock When self-reference fails to close: Matrix-level dynamics in large language models.
\newblock \texttt{arXiv:2604.12128}, 2026.

\bibitem{berg2025}
C.~Berg, D.~de~Lucena, and J.~Rosenblatt.
\newblock Large language models report subjective experience under self-referential processing.
\newblock \texttt{arXiv:2510.24797}, 2025.

\bibitem{dadfar2026}
Z.~P.~Dadfar.
\newblock When models examine themselves: Vocabulary-activation correspondence in self-referential processing.
\newblock \texttt{arXiv:2602.11358}, 2026.

\bibitem{thrush2024}
T.~Thrush et~al.
\newblock {I} am a strange dataset: Metalinguistic tests for language models.
\newblock In \emph{Proc. ACL}, 2024.

\bibitem{merrill2023}
W.~Merrill and A.~Sabharwal.
\newblock The parallelism tradeoff: Limitations of log-precision transformers.
\newblock \emph{Trans. ACL}, 11:531--545, 2023.

\end{thebibliography}

\end{document}
